\title{SAINT: Improved Neural Networks for Tabular Data via Row Attention and Contrastive Pre-Training}

\begin{abstract}
Tabular data underpins numerous high-impact applications of machine learning from fraud detection to genomics and healthcare.  Classical approaches to solving tabular problems, such as gradient boosting and random forests, are widely used by practitioners.  However, recent deep learning methods have achieved a degree of performance competitive with popular techniques.  We devise a hybrid deep learning approach to solving tabular data problems.  Our method, SAINT, performs attention over both rows and columns, and it includes an enhanced embedding method.  We also study a new contrastive self-supervised pre-training method for use when labels are scarce.  SAINT consistently improves performance over previous deep learning methods, and it even outperforms gradient boosting methods, including XGBoost, CatBoost, and LightGBM, on average over a variety of benchmark tasks.
\end{abstract}

\section{Introduction}
\label{sec:introduction}
While machine learning for image and language processing has seen major advances over the last decade, many critical industries, including financial services, health care, and logistics, rely heavily on data in structured table format. 

Tabular data is unique in several ways that have prevented it from benefiting from the impressive success of deep learning in vision and language. First, tabular data often contain heterogeneous features that represent a mixture of continuous, categorical, and ordinal values, and these values can be independent or correlated. Second, there is no inherent positional information in tabular data, meaning that the order of columns is arbitrary. This differs from text, where tokens are always discrete, and ordering impacts semantic meaning.  It also differs from images, where pixels are typically continuous, and nearby pixels are correlated. Tabular models must handle features from multiple discrete and continuous distributions, and they must discover correlations without relying on the positional information. Sufficiently powerful deep learning systems for tabular data have the potential to improve performance beyond what is achieved by classical methods, like linear classifiers and random forests.  Furthermore, without performant deep learning models for tabular data, we lack the ability to exploit compositionality, end-to-end multi-task models, fusion with multiple modalities (e.g. image and text), and representation learning. 

We introduce SAINT, the Self-Attention and Intersample Attention Transformer, a specialized architecture for learning with tabular data. SAINT leverages several mechanisms to overcome the difficulties of training on tabular data. SAINT projects all features -- categorical and continuous -- into a combined dense vector space. These projected values are passed as tokens into a transformer encoder which uses attention in the following two ways. First, there is ``self-attention,'' which attends to individual features within each data sample. Second, we propose a novel ``intersample attention,'' which enhances the classification of a row (i.e., a data sample) by relating it to other rows in the table.  Intersample attention is akin to a nearest-neighbor classification, where the distance metric is learned end-to-end rather than fixed. In addition to this hybrid attention mechanism, we also leverage self-supervised contrastive pre-training to boost performance for semi-supervised problems.  

We provide comparisons of SAINT to a wide variety of deep tabular architectures and commonly used tree-based methods using a diverse battery of datasets. We observe that SAINT, on average, outperforms all other methods on supervised and semi-supervised tasks.  More importantly, SAINT often out-performs boosted trees (including  XGBoost \cite{chen2016xgboost}, CatBoost \cite{dorogush2018catboost}, and LightGBM \cite{ke2017lightgbm}), which have long been an industry favorite for complex tabular datasets.

Finally, we visualize the attention matrices produced by our models to gain insights into how they behave.

\section{Related Work}
\label{sec:related_work}
\textbf{Classical Models} The most widely adopted approaches for supervised and semi-supervised learning on tabular datasets eschew neural models due to their black-box nature and high compute requirements. When one has reasonable expectations of linear relationships, a variety of modeling approaches are available \cite{wright1995logistic,weisberg2005applied,starkweather2011multinomial,mcculloch2005generalized}. In more complex settings, non-parametric tree-based models are used. Commonly used tools such as XGBoost~\cite{chen2016xgboost}, CatBoost~\cite{dorogush2018catboost}, and LightGBM~\cite{ke2017lightgbm} provide several benefits such as interpretability, the ability to handle a variety of feature types including null values, as well as performance in both high and low data regimes.

\textbf{Deep Tabular Models} While classical methods are still the industry favorite, some recent work brings deep learning to the tabular domain. For example, TabNet \cite{arik2019TabNet} uses neural networks to mimic decision trees by placing importance on only a few features at each layer. The attention layers in that model do not use the regular dot-product self-attention common in transformer-based models, rather there is a type of sparse layer that allows only certain features to pass through. \citet{yoon2020vime} propose VIME, which employs MLPs in a technique for pre-training based on denoising. TABERT \cite{yin2020tabert}, a more elaborate neural approach inspired by the large language transformer model BERT \cite{devlin2018bert}, is trained on semi-structured test data to perform language-specific tasks. Several other studies utilize tabular data, but their problem settings are outside of our scope \cite{pathak2016context,chen2019learning,padhi2021tabular,shavitt2018regularization, katzir2021netdnf}.

Transformer models for more general tabular data include TabTransformer~\cite{huang2020tabtransformer}, which uses a transformer encoder to learn contextual embeddings \emph{only} on categorical features. The continuous features are concatenated to the embedded features and fed to an MLP. The main issue with this model is that continuous data do not go through the self-attention block. That means any information about correlations between categorical and continuous features is lost. In our model, we address that issue by projecting continuous features and categorical features to the higher dimensional embedding space and passing them both through the transformer blocks. In addition, we propose a new type of attention to explicitly allow data points to attend to each other to get better representations.

\textbf{Axial Attention} \citet{ho2019axial} are the first to propose row and column attention in the context of localized attention in 2-dimensional inputs (like images) in their Axial Transformer. This is where for a given pixel, the attention is computed only on the pixels that are on the same row and column, rather than using all the pixels in the image. The MSA Transformer~\cite{rao2021msa} extends this work to protein sequences and applies both column and row attention across similar rows (tied row attention). TABBIE~\cite{iida2021tabbie} is an adaptation of axial attention that applies self-attention to rows and columns separately, then averages the representations and passes them as input to the next layer. In all these works, different features from the same data point communicate with each other and with the same feature from a whole batch of data. Our approach, intersample attention, is hierarchical in nature; first features of a given data point interact with each other, then data points interact with each other using entire rows/samples. 

In a similar vein, Graph Attention Networks~(GAT)~\cite{velivckovic2017graph} seek to compute attention over neighbors on a graph, thereby learning which neighbor's information is most relevant to a given node’s prediction. One way to view our intersample attention is as a GAT on a complete graph where all tabular rows are connected to all other rows. \citet{yang2016hierarchical} explore hierarchical attention for the task of document classification where attention is computed between words in a given sentence and then between the sentences, but they did not attempt to compute the attention between entire documents themselves.

\textbf{Self-Supervised Learning} Self-supervision via a `pretext task' on unlabeled data coupled with finetuning on labeled data is widely used for improving model performance in language and computer vision.  

Some of the tasks previously used for self-supervision on tabular data include masking, denoising, and replaced token detection. Masking (or Masked Language Modeling(MLM)) is when individual features are masked and the model's objective is to impute their value~\cite{pathak2016context,arik2019TabNet,huang2020tabtransformer}. Denoising injects various types of noise into the data, and the objective there is to recover the original values~\cite{vincent2008extracting,yoon2020vime}. Replaced token detection (RTD) inserts random values into a given feature vector and seeks to detect the location of these replacements~\cite{huang2020tabtransformer,iida2021tabbie}. Contrastive pre-training, where the distance between two views of the same point is minimized while maximizing the distance between two different points~\cite{chen2020simple,he2020momentum,grill2020bootstrap}, is another pretext task that applies to tabular data.  In this paper, to the best of our knowledge, we are the first to adopt contrastive learning for tabular data.  We couple this strategy with denoising to perform pre-training on a plethora of datasets with varied volumes of labeled data, and we show that our method outperforms traditional boosting methods.

\section{Self-Attention and Intersample Attention Transformer (SAINT)}
\label{sec:saint}

In this section, we introduce our model, Self-Attention and Intersample Attention Transformer~(SAINT), and explain in detail its various components.

Suppose $\mathcal{D} = \{\mathbf{x_i},y_i\}_{i=1}^m$ is a tabular dataset with $m$ points, where each ${x_i}$ is an $n$-dimensional feature vector, and ${y_i}$ is a label or target variable. Similar to BERT~\citep{devlin2018bert}, we append a \verb|[CLS]| token with a learned embedding to each data sample. Let $\mathbf{x_i} = [\verb|[CLS]|, f_i^{\{1\}}, f_i^{\{2\}},  ..,f_i^{\{n\}}]$ be a single data-point with categorical or continuous features $f_i^{\{j\}}$, and let $\mathbf{E}$ be the embedding layer that embeds each feature into a $d$-dimensional space.  Note that $\mathbf{E}$ may use different embedding functions for different features. For a given $\mathbf{x_i} \in \mathbb{R}^{(n+1)}$, we get $\mathbf{E}(\mathbf{x_i}) \in \mathbb{R}^{(n+1) \times d}$.

\textbf{Encoding the Data} In language models, all tokens are embedded using the same procedure. However, in the tabular domain, different features can come from distinct distributions, necessitating a heterogeneous embedding approach.  Note that tabular data can contain multiple categorical features which may use the same set of tokens. Unless it is known that common tokens possess identical relationships within multiple columns, it is important to embed these columns independently.  Unlike the embedding of TabTransformer\cite{huang2020tabtransformer}, which uses attention to embed only categorical features, we propose also projecting continuous features into a $d-$dimensional space before passing their embedding through the transformer encoder. To this end, we use a separate single fully-connected layer with a ReLU nonlinearity for each continuous feature, thus projecting the $1-$dimensional input into $d-$dimensional space. With this simple trick alone, we significantly improve the performance of the TabTransformer model as discussed in Section \ref{subsec:main_results}. An additional discussion concerning positional encodings can be found in Appendix \ref{appendix_sec:complete_training_deets}.

\subsection{Architecture}

SAINT is inspired by the transformer encoder of \citet{vaswani2017attention}, designed for natural language, where the model takes in a sequence of feature embeddings and outputs contextual representations of the same dimension. A graphical overview of SAINT is presented in Figure \ref{fig:saint_arch_training}(a).

SAINT is composed of a stack of $L$ identical stages. Each stage consists of one self-attention transformer block and one intersample attention transformer block. The self-attention transformer block is identical to the encoder from \cite{vaswani2017attention}. It has a multi-head self-attention layer~(MSA) (with $h$ heads), followed by two fully-connected feed-forward (FF) layers with a GELU non-linearity~\cite{hendrycks2016gaussian}. Each layer has a skip connection~\cite{he2016deep} and layer normalization (LN)~\cite{ba2016layer}. The intersample attention transformer block is similar to the self-attention transformer block, except that the self-attention layer is replaced by an intersample attention layer~(MISA). The details of the intersample attention layer are presented in the following subsection.

The SAINT pipeline, with a single stage ($L=1$) and a batch of $b$ inputs, is described by the following equations. We denote multi-head self-attention by MSA, multi-head intersample attention by MISA, feed-forward layers by FF, and layer norm by LN:

\begin{align}
    \mathbf{z_i^{(1)}} &= \operatorname{LN}(\operatorname{MSA}(\mathbf{E}(\mathbf{x_i}))) + \mathbf{E}(\mathbf{x_i}) & \mathbf{z_i^{(2)}} &= \operatorname{LN}(\operatorname{FF_1}(\mathbf{z_i^{(1)}})) +\mathbf{z_i^{(1)}}
    \label{eq:sa_block_apply} \\
    \mathbf{z_i^{(3)}} &= \operatorname{LN}(\operatorname{MISA}(\{\mathbf{z_i^{(2)}}\}_{i=1}^b)) +\mathbf{z_i^{(2)}} &  \mathbf{r_i}&= \operatorname{LN}(\operatorname{FF_2}(\mathbf{z_i^{(3)}})) +\mathbf{z_i^{(3)}}
    \label{eq:isa_block_apply} 

\end{align}

where $\mathbf{r_i}$ is SAINT's contextual representation output corresponding to data point $\mathbf{x_i}$. This contextual embedding can be used in downstream tasks such as self-supervision or classification.

\subsection{Intersample attention}
\label{subsec:ISA}

We introduce intersample attention (a type of row attention) where the attention is computed across different data points (rows of a tabular data matrix) in a given batch rather than just the features of a single data point. Specifically, we concatenated the embeddings of each feature for a single data point, then compute attention over samples (rather than features). This enables us to improve the representation of a given point by inspecting other points. 

When a feature is missing or noisy in one row, intersample attention enables SAINT to borrow the corresponding features from other similar data samples in the batch.

An illustration of how intersample attention is performed in a single head is shown in Figure ~\ref{fig:ISA} and the pseudo-code is presented in Algorithm~\ref{alg:MISA}.  Unlike the row attention used in \cite{ho2019axial,child2019generating,rao2021msa,iida2021tabbie}, intersample attention allows all features from different samples to communicate with each other.  In our experiments, we show that this ability boosts performance appreciably.  In the multi-head case, instead of projecting $q,k,v$ to a given dimension $d$, we project them to $d/h$ where $h$ is the number of heads. Then we concatenate all the updated value vectors, $v_i$, to get back a vector of length $d$.

\begin{algorithm}[tb]
\captionsetup{font=small}
   \caption{PyTorch-style pseudo-code for intersample attention. For simplicity, we describe just one head and assume the value vector dimension is same as the input embedding dimension.}
   \label{alg:MISA}

    \definecolor{codeblue}{rgb}{0.25,0.5,0.5}
    \lstset{
      basicstyle=\fontsize{8.5 pt}{8.5 pt}\ttfamily\bfseries,
      commentstyle=\fontsize{8.5 pt}{8.5 pt}\color{codeblue},
      keywordstyle=\fontsize{8.5 pt}{8.5 pt}\color{magenta},
    }
    \begin{lstlisting}[language=python]
    # b: batch size, n: number of features, d: embedding dimension # W_q,  W_k, W_v are weight matrices of dimension dxd # mm: matrix-matrix multiplication def self_attention(x): # x is bxnxd q, k, v = mm(W_q,x), mm(W_k,x), mm(W_v,x) #q,k,v are bxnxd attn = softmax(mm(q,np.transpose(k, (0, 2, 1)))/sqrt(d)) # bxnxn out = mm(attn, v) #out is bxnxd return out def intersample_attention(x): # x is bxnxd b,n,d = x.shape # as mentioned above x = reshape(x, (1,b,n*d)) # reshape x to 1xbx(n*d) x = self_attention(x) # the output x is 1xbx(n*d) out = reshape(x,(b,n,d)) # out is bxnxd return out
    \end{lstlisting}
\end{algorithm}

\section{Pre-training \& Finetuning}
\label{sec:pretraining}

Contrastive learning, in which models are pre-trained to be invariant to reordering, cropping, or other label-preserving ``views'' of the data~ \cite{chen2020simple,he2020momentum,pathak2016context,grill2020bootstrap,vincent2008extracting}, is a powerful tool in the vision and language domains that has never (to our knowledge) been applied to tabular data. We present a contrastive pipeline for tabular data, a visual description of which is shown in Figure \ref{fig:saint_arch_training}.

Existing self-supervised objectives for tabular data include denoising \citep{vincent2008extracting}, a variation of which was used by VIME~\citep{yoon2020vime}, masking, and replaced token detection as used by TabTransformer \citep{huang2020tabtransformer}. We find that, while these methods are effective, superior results are achieved by contrastive learning.

\textbf{Generating augmentations} Standard contrastive methods in vision craft different ``views'' of images using crops and flips. It is difficult to craft invariance transforms for tabular data. The authors of VIME~\citep{yoon2020vime} use mixup in the non-embedded space as a data augmentation method, but this is limited to continuous data. We instead use CutMix~\citep{yun2019cutmix} to augment samples in the input space and we use mixup~\citep{zhang2017mixup} in the embedding space. These two augmentations combined yield a challenging and effective self-supervision task. Assume that only $l$ of $m$ data points are labeled.  We denote the embedding layer by $\mathbf{E}$, the SAINT network by $\mathbf{S}$, and 2 projection heads as $g_1(\cdot)$ and $g_2(\cdot)$. The CutMix augmentation probability is denoted $p_{\text{cutmix}}$ and the mixup parameter is $\alpha$.
Given point $\mathbf{x_i}$, the original embedding is $\mathbf{p_i} = \mathbf{E}(\mathbf{x_i})$, while the augmented representation is generated as follows:

\begin{align}
\mathbf{x_i^\prime} &= \mathbf{x_i}\odot \mathbf{m} + \mathbf{x_a}\odot \mathbf{(1-m)} && \textcolor{gray}{\text{CutMix in raw data space}}\\
\mathbf{p_i^\prime} &= \alpha*\mathbf{E}(\mathbf{x_i^\prime}) +(1-\alpha)*\mathbf{E}(\mathbf{x_b^\prime})&& \textcolor{gray}{\text{mixup in embedding space}}
\end{align}

where $\mathbf{x_a}$, $\mathbf{x_b}$ are random samples from the current batch, $\mathbf{x_b^\prime}$ is the CutMix version of $\mathbf{x_b}$, $\mathbf{m}$ is the binary mask vector sampled from a Bernoulli distribution with probability $p_{\text{cutmix}}$, and $\alpha$ is the mixup parameter. Note that we first obtain a CutMix version of every data point in a batch by randomly selecting a partner to mix with. We then embed the samples and choose new partners before performing mixup. 

\textbf{SAINT and projection heads} Now that we have both the clean~$\mathbf{p_i}$ and mixed~$\mathbf{p_i^\prime}$ embeddings, we pass them through SAINT, then through two projection heads, each consisting of an MLP with one hidden layer and a ReLU. The use of a projection head to reduce dimensionality before computing contrastive loss is common in vision~\citep{chen2020simple,he2020momentum,grill2020bootstrap} and indeed also improves results on tabular data. Ablation studies and further discussion are available in Appendix~\ref{appendix_sec:analysis}.

\textbf{Loss functions} We consider two losses for the pre-training phase. (i) The first is a contrastive loss that pushes the latent representations of two views of the same data point ($z_i$ and $z_i^\prime$) close together and encourages different points ($z_i$ and $z_j$, $i\neq j$) to lie far apart.  For this, we borrow the InfoNCE loss from metric-learning works~\citep{sohn2016improved,oord2018representation,chen2020simple, wu2018unsupervised}; (ii) The second loss comes from a denoising task. For denoising, we try to predict the original data sample from a noisy view. Formally, we are given $\mathbf{r_i^\prime}$ and we reconstruct the inputs as $\mathbf{x_i^{\prime \prime}}$ to minimize the difference between the original and the reconstruction. The combined pre-training loss is:

\begin{align}

\mathcal{L_\text{pre-training}} &= \underbrace{-\sum_{i=1}^{m}{\log{\frac{\exp(z_i \cdot z_i^\prime/\tau)}{\sum_{k=1}^{m}{\exp(z_i \cdot z_k^\prime/\tau)}}}}}_{\text{Contrastive Loss}} + \lambda_{\text{pt}} \underbrace{\sum_{i=1}^{m}\sum_{j=1}^{n} [\mathcal{L}_j(\text{MLP}_j(\mathbf{r_i^\prime}),\mathbf{x_i})]}_{\text{Denoising Loss}}
\end{align}

where $\mathbf{r_i} = \mathbf{S}(\mathbf{p_i}), \mathbf{r_i^\prime} = \mathbf{S}(\mathbf{p_i^\prime}), z_i = g_1(\mathbf{r_i}),z_i^\prime = g_2(\mathbf{r_i^\prime}) $. $\mathcal{L}_j$ is cross-entropy loss or mean squared error depending on the $j^{th}$ feature being categorical or continuous. Each $\text{MLP}_j$ is a single hidden layer perceptron with a ReLU non-linearity. There are $n$ in number, one for each input feature. $\lambda_{\text{pt}}$ is a hyper-parameter and $\tau$ is temperature parameter and both of these are tuned using validation data.

\textbf{Finetuning} Once SAINT is pre-trained on all unlabeled data, we finetune the model on the target prediction task using the $l$ labeled samples.  The pipeline of this step is shown in Figure~\ref{fig:saint_arch_training}(b). For a given point $\mathbf{x_i}$, we learn the contextual embedding $\mathbf{r_i}$. For the final prediction step, we pass the embedding corresponding only to the \verb|[CLS]| token through a simple MLP with a single hidden layer with ReLU activation to get the final output. We evaluate cross-entropy loss on the outputs for classification tasks and mean squared error for regression tasks.

\section{Experimental Evaluation}
\label{sec:expts}
We evaluate SAINT on 16 tabular datasets. In this section, we discuss variants of SAINT and evaluate them in both supervised and semi-supervised scenarios. We also analyze each component of SAINT and perform ablation studies to understand the importance of each component in the model. Using visualization, we interpret the behavior of attention maps.   

\textbf{Datasets} We evaluate SAINT on 14 binary classification tasks and 2 multi-class classification tasks. These datasets were chosen because they were previously used to evaluate competing methods \cite{yoon2020vime,arik2019TabNet,huang2020tabtransformer}. They are also diverse; the datasets range in size from 200 to 495,141 samples, and from 8 to 784 features, with both categorical and continuous features. Some datasets are missing data while some are complete and some are well-balanced while others have highly skewed class distributions. Each of these datasets is publicly available from either UCI\footnote{http://archive.ics.uci.edu/ml/datasets.php} or AutoML.\footnote{https://automl.chalearn.org/data} Details of these datasets can be found in Appendix \ref{appendix_sec:datasets}. As the pre-processing step for each dataset, all the continuous features are Z-normalized, and all categorical features are label-encoded before the data is passed on to the embedding layer.

\textbf{Model variants} The SAINT architecture discussed in the previous section has one self-attention transformer encoder block stacked with one intersample attention transformer encoder block in each stage. We also consider variants with just one of these blocks. SAINT-s variant has only self-attention, and SAINT-i has only intersample attention. SAINT-s is exactly the encoder from \cite{vaswani2017attention} but applied to tabular data. See Table \ref{table:SAINT_config} for an architectural comparison of these model variants.

\textbf{Baselines} We compare our model to traditional methods like logistic regression and random forests. We benchmark against the powerful boosting libraries XGBoost, LightGBM, and CatBoost. We also compare against deep learning methods, like multi-layer perceptrons, VIME, TabNet, and TabTransformer. For the methods that use unsupervised pre-training as a preprocessing step, we used Masked Language Modeling~(MLM) for TabNet~\cite{devlin2018bert}, and we use Replaced Token Detection~(RTD) for TabTransformer~\cite{clark2020electra} as mentioned in the respective papers. For multi-layer perceptrons, we use denoising~\cite{vincent2008extracting} as suggested in VIME. 

\textbf{Metrics} Since the majority of the tasks used in our analysis are binary classification, we use AUROC as the primary metric to measure performance. AUROC captures how well the model separates the two classes in the dataset. For the two multi-class datasets, Volkert and MNIST, we use the accuracy on the test set to compare performance.

\textbf{Training} We train all the models (including pre-training runs) using AdamW with $\beta_1 = 0.9$, $\beta_2 = 0.999$, $\text{decay} = 0.01$, and with a learning rate of $0.0001$ with batches of size 256 (except for datasets with a large number of columns like MNIST and Arcene, for which we use smaller batch sizes). We split the data into $65\%$, $15\%$, and $25\%$ for training, validation, and test splits, respectively. We vary the embedding size based on the number of features in the dataset. The exact configurations for each of the datasets are presented in Appendix \ref{appendix_sec:complete_training_deets}. We use CutMix mask parameter $p_{\text{cutmix}} = 0.3$ and mixup parameter $\alpha=0.2$ for all the standard pre-training experiments. We use pre-training loss hyper-parameters $\lambda_{\text{pt}} = 10$ and temperature $\tau = 0.7$ for all settings. 

\subsection{Results}
\label{subsec:main_results}

\textbf{Supervised setting} In Table~\ref{table:supervised_results}, we report results on 7 representative binary classification and 2 multi-class classification datasets, as well as the average performance across all 14 binary classification datasets. Note that each number reported in the Table~\ref{table:supervised_results} is the mean from 5 trials with different seeds. In 13 out of 16 datasets, one of the SAINT variants outperforms all baseline models. In the remaining 3 datasets, our model's performance is very close to the best model. On average, SAINT variants each perform better than all baseline models by a significant margin, and SAINT performs even better than its two variants. For complete results from every dataset as well as details including standard error, see Appendix \ref{appendix_sec:exhaustive_results}.

\textbf{Semi-supervised setting}  We perform 3 sets of experiments with 50, 200, and 500 labeled data points (in each case the rest are unlabeled). See Table~\ref{table:semisupervised} for numerical results. In all cases, the pre-trained SAINT model (with both self and intersample attention) performs the best. Interestingly, we note that when all the training data samples are labeled, pre-training does not contribute appreciably, hence the results with and without pre-training are fairly close. 

\textbf{Effect of embedding continuous features} To understand the effect of learning embeddings for continuous data, we perform a simple experiment with TabTransformer. We modify TabTransformer by embedding continuous features into $d$ dimensions using a single layer ReLU MLP, just as they use on categorical features, and we pass the embedded features through the transformer block. We keep the entire architecture and all training hyper-parameters the same for both TabTransformer and its modified version. The average AUROC of the original TabTransformer is 89.38. Just by embedding the continuous features, the performance jumps to 91.72. This experiment shows that embedding the continuous data is important and can boost the performance of the model significantly. 

\textbf{When to use intersample attention?} From our experiments, we observe that SAINT-i consistently outperforms other variants whenever the number of features is large. In particular, whenever there are few training data points coupled with many features (which is common in biological datasets), SAINT-i outperforms SAINT-s significantly (see the ``Arcene'' and ``Arrhythmia'' results). Another advantage of using SAINT-i is that execution is fast compared to SAINT-s,  despite the fact that the number of parameters of SAINT-i is much higher than that of SAINT-s (see Table~\ref{table:SAINT_config}). 

\textbf{How robust is SAINT to data corruptions?} We evaluate the robustness of SAINT variants by corrupting the training data.  To simulate corruption, we apply CutMix, replacing 10\% to 90\% of the features with values of other randomly selected samples. The drop in the mean AUROC is quite minimal until 70\% data corruption when the performance drops significantly. SAINT and SAINT-i models are comparatively more robust than SAINT-s. This shows that using row attention improves the model's robustness to noisy training data as we anticipated. However, we find the opposite trend when many features are missing in the training data. SAINT-s and SAINT are quite robust, and the drop in AUROC is not drastic even when 90\% of the data is missing. This observation shows that SAINT is reliable for training on corrupted training data. The AUROC trend line plots for both the scenarios are shared in Appendix \ref{appendix_sec:analysis}.

\textbf{Effect of batch size on intersample attention performance} As discussed in Section~\ref{subsec:ISA}, attention is computed between batches of data points. We examine the impact of batch size using batches of size ranging from 32 to 256. We find that the variation in SAINT-i's performance is low and is comparable to that of SAINT-s, which has no intersample attention component. We present the plots in Appendix \ref{appendix_sec:analysis}.

\subsection{Interpreting attention}
\label{subsec:attention_analysis}
One advantage of using transformer-based models is that attention comes with some interpretability, in contrast, MLPs are hard to interpret. In particular, when we use only one transformer stage, the attention maps reveal which features and which data points are being used by the model to make decisions. We use MNIST data to examine how self-attention and intersample attention behave in our models. While MNIST is not a typical tabular dataset, it has the advantage that its features can be easily visualized as an image.

Figure~\ref{fig:SA_in_SAINT} depicts the attention on each of the pixels/features in a self-attention layer of SAINT. Without any explicit supervision, the model learns to focus on the foreground pixels, and we clearly see from the attention map which features are most important to the model. The self-attention plots of SAINT-s are similar (Appendix \ref{appendix_sec:additional_interpret}).

Figures~\ref{fig:ISA_in_SAINT} and \ref{fig:ISA_in_SAINTi} depict a similar visualization on a batch of 20 points, 2 from each class in MNIST. Figure~\ref{fig:ISA_in_SAINT} shows intersample attention in SAINT. This plot shows which samples attend to which other samples in the batch. Surprisingly, very few points in a batch receive attention. We hypothesize that the model focuses on a few points that are critical because they are particularly difficult to classify without making direct comparisons to exemplars in the batch. In Figure~\ref{fig:ISA_in_SAINTi}, we show the intersample attention plot from a SAINT-i model. The same sparse attention behaviour persists here too, but the points being attended to are different in this model. Interestingly, we find this behavior to be significantly different on the Volkert data, where a wide range of data becomes the focus of attention depending on the input. The intersample attention layer gets dense with the hardness~(to classify) of the datasets. See Appendix ~\ref{appendix_sec:additional_interpret} for additional MNIST and Volkert attention maps.

Figure~\ref{fig:attention_tsne} shows the behavior of attention at the dataset (rather than batch) level. We visualize a t-SNE \citep{van2008visualizing} embedding for \textit{value} vectors generated in intersample attention layers, and we highlight the points that are most attended to in each batch. In Figure~\ref{fig:attention_tsne}~(left), the \textit{value} vectors and attention are computed on the output representations of a self-attention layer. 

In contrast, the \textit{value} vectors and attention in Figure~\ref{fig:attention_tsne}~(right) are computed on the embedding layer output, since the SAINT-i model does not use self-attention.  In these two plots, the classes to which the model attends vary dramatically. Thus, the exact classes to which an attention head pays attention change with the architecture, but the trend of using a few classes as a `pivot' seems to be prevalent in intersample attention heads. Additional analyses are presented in Appendix \ref{appendix_sec:analysis}.

\section{Discussion, Limitations, and Impact}
\label{sec:discussion}
We introduce intersample attention, contrastive pre-training, and an improved embedding strategy for tabular data. Even though tabular data is an extremely common data format used by institutions in various domains, deep learning methods are still lagging behind tree-based boosting methods in production. With SAINT, we show that neural models can often improve upon the performance of boosting methods across numerous datasets with varying characteristics.

SAINT offers improvements in a widely used domain, which is quite impactful. While our method performs well on the diverse tabular datasets studied here, real-world applications contain a broad range of datasets which may be highly noisy or imbalanced.  Moreover, we have tuned SAINT for the settings in which we test it.  Thus, we caution practitioners against assuming that what works on the benchmarks in this paper will work in their own setting.

{
\small

\end{document}
